\section{Problem and Metrics}
\label{sec:motivation}

Consider an iteration with outer token budget $B$, Prefill bound $C$, and $D_t$
active Decode tokens. $B$ is the total scheduled-token budget; $C$ additionally
caps prompt work when configured. A decode-first mixed scheduler admits
\begin{equation}
    P_t \leq \min(C, B-D_t),
    \label{eq:mixed-budget}
\end{equation}
where $P_t$ is the number of prompt tokens advanced in the iteration. The
bound $C$ limits a single chunk, but the current Decode batch determines how
much of the token budget remains for Prefill. The two decisions are therefore
coupled: the scheduler chooses $(P_t,D_t)$ inside an iteration and the queue
state indirectly determines the phase share over time.

PD-TDM exposes the second decision. It first selects a phase
$u_t\in\{P,D\}$ and then builds its batch subject to
\begin{equation}
 u_t=P \Rightarrow D_t=0, \qquad
 u_t=D \Rightarrow P_t=0, \qquad P_t\leq C.
 \label{eq:phase-invariant}
\end{equation}
The phase decision is made only when both queues compete. A target
$r_k\in[0,1]$ expresses the desired Prefill tendency in controller window $k$;
it is a credit rate, not a wall-clock P:D ratio, because Prefill and Decode
iterations have different durations and queue fallbacks can override it. The
evaluated controller updates this target at a slower timescale from windowed
TTFT and TPOT pressure, while the bucket applies it at every iteration. Let
$v_{P,k}$ and $v_{D,k}$ be the measured violation rates and let $q$ be the
target violation rate. The update is summarized by
\begin{equation}
 r_{k+1}=\Pi_{[r_{\min},r_{\max}]}
 \left((1-\alpha)r_k+\alpha\left[r_k+\kappa(e_{P,k}-e_{D,k})\right]\right),
 \label{eq:ratio-update}
\end{equation}
where $e_{P,k}=[v_{P,k}-q]_+$ and $e_{D,k}=[v_{D,k}-q]_+$, and $\Pi$ clips to
the configured range. This is the active update: persistent TPOT saturation
zeros its error; cold or KV-frozen windows hold $r_k$, as does a sub-deadband
change. The main runs disable the optional queue-age term and use $q=0.05$,
$\kappa=0.5$, $\alpha=0.3$,
$r_0=0.30$, $r_{\min}=0.05$, and $r_{\max}=0.80$; the controller updates
every eight scheduler iterations. A $0.02$ deadband, a 16-sample warm-up,
and a two-tick TPOT-saturation detector are also enabled. We evaluate the
complete two-timescale path, but do not isolate the marginal gain of this
update from phase purity, chunking, or backend execution.

For all requests arriving in measurement window $T$, let $o_i$ be the generated
output length and let $S_T,S_D$ be TTFT and TPOT SLOs. Failed requests and
requests without matched latency telemetry are counted as non-attainments and
contribute zero to goodput.
We report request attainment and output-token goodput as
\begin{equation}
 \begin{aligned}
 A&=\frac{1}{N_T}\sum_i I_i,
 &G&=\frac{1}{T}\sum_i o_i I_i,\\[-1pt]
 I_i&=\mathbf{1}[\mathrm{TTFT}_i<S_T\land\mathrm{TPOT}_i<S_D].
 \end{aligned}
 \label{eq:metrics}
\end{equation}
TPOT is the per-request mean inter-token time; its cross-request percentile is
not a guarantee on every individual token interval. For T6, $N_T$ counts
arrivals in $[30,90)$~s and $T=60$~s; latency summaries are computed on
successful records with matched telemetry. Thus $G$ is arrival-cohort goodput,
not completion-window throughput: output tokens from a qualifying cohort
request are counted even if it finishes after the cohort end. Equation~\ref{eq:metrics} avoids
replacing token-weighted goodput with raw throughput multiplied by request
attainment, which is not valid for variable-length outputs.
